\documentclass[10pt]{article}

% ---------- Document Identity (update these only) ----------
\newcommand{\DocStage}{}
\newcommand{\DocTitle}{Your Road to Tier One SOF}
\newcommand{\ProgramName}{182-Week Civilian-to-Selection Readiness Pipeline}
\newcommand{\ProgramSubtitle}{}

% ---------- Page ----------
\usepackage[legalpaper,left=0.5in,right=0.5in,top=0.6in,bottom=0.45in]{geometry}

% ---------- Typography ----------
\usepackage[T1]{fontenc}
\usepackage[utf8]{inputenc}
\usepackage{libertinus}
\usepackage{microtype}
\usepackage{parskip}
\usepackage{ragged2e}
\usepackage{textcomp}
\AtBeginDocument{\color{black}}
\AtBeginDocument{\pagecolor{white}}
\AtBeginDocument{\justifying}

% ---------- Landscape pages ----------
\usepackage{pdflscape}

% ---------- Color ----------
\usepackage[table]{xcolor}
\definecolor{StageBlue}{HTML}{1F4E79}
\definecolor{StageBlueDark}{HTML}{163A5A}
\definecolor{LightGray}{HTML}{F2F4F7}
\definecolor{MidGray}{HTML}{D9DEE5}
\definecolor{RuleGray}{HTML}{C7CED8}
\definecolor{HeaderInk}{HTML}{000000}
\definecolor{Ink}{HTML}{000000}

% ---------- Utility ----------
\usepackage{enumitem}
\sloppy
\setlength{\emergencystretch}{2em}

% ---------- Boxes / UI ----------
\usepackage[most]{tcolorbox}

\tcbset{
  charterTitle/.style={
    enhanced,
    colback=StageBlue!4,
    colframe=StageBlue,
    boxrule=1.25pt,
    arc=3pt,
    left=16pt,right=16pt,top=14pt,bottom=14pt,
    borderline north={3.2pt}{0pt}{StageBlue},
    borderline west={4pt}{0pt}{StageBlue},
    borderline south={1.25pt}{0pt}{StageBlue},
    drop shadow={black!25!white},
  },
  charterRule/.style={
    enhanced,
    colback=LightGray,
    colframe=RuleGray,
    boxrule=0.85pt,
    arc=3pt,
    left=12pt,right=12pt,top=10pt,bottom=10pt,
    borderline west={3pt}{0pt}{StageBlue},
  },
  charterBadge/.style={
    enhanced,
    colback=StageBlue,
    coltext=white,
    colframe=StageBlueDark,
    boxrule=0.6pt,
    arc=2pt,
    left=6pt,right=6pt,top=3pt,bottom=3pt,
  }
}
\newtcbox{\Badge}{nobeforeafter,charterBadge}

% ---------- Lists ----------
\setlist[itemize]{leftmargin=1.5em, itemsep=0.25em, topsep=0.25em}
\setlist[enumerate]{leftmargin=1.8em, itemsep=0.25em, topsep=0.25em}

% ---------- TOC ----------
\usepackage{tocloft}
\setcounter{tocdepth}{3}
\setlength{\cftbeforesecskip}{3pt}
\setlength{\cftbeforesubsecskip}{2pt}
\setlength{\cftbeforesubsubsecskip}{0pt}
\renewcommand{\cftsecfont}{\bfseries}
\renewcommand{\cftsecpagefont}{\bfseries}
\renewcommand{\cftsubsecfont}{\normalfont}
\renewcommand{\cftsubsecpagefont}{\normalfont}
\renewcommand{\cftsubsubsecfont}{\normalfont}
\renewcommand{\cftsubsubsecpagefont}{\normalfont}
\setlength{\cftsecindent}{0pt}
\setlength{\cftsubsecindent}{1.25em}
\setlength{\cftsubsubsecindent}{2.8em}
\setlength{\cftsecnumwidth}{2.1em}
\setlength{\cftsubsecnumwidth}{2.8em}
\setlength{\cftsubsubsecnumwidth}{3.6em}

% Remove the built-in TOC heading so only the custom heading is shown
\makeatletter
\renewcommand{\tableofcontents}{\@starttoc{toc}}
\makeatother

% ---------- Header/Footer: page number only ----------
\usepackage{fancyhdr}
\pagestyle{fancy}
\fancyhf{}
\renewcommand{\headrulewidth}{0pt}
\renewcommand{\footrulewidth}{0pt}
\fancyfoot[C]{\thepage}

% ---------- Section formatting ----------
\usepackage{titlesec}

\titleformat{\section}
  {\Large\bfseries\color{Ink}}
  {\thesection}{0.6em}{}
  [\vspace{0.25em}\textcolor{StageBlue}{\rule{\linewidth}{0.8pt}}]

\titleformat{\subsection}
  {\large\bfseries\color{Ink}}
  {\thesubsection}{0.6em}{}

\titleformat{\subsubsection}
  {\normalsize\bfseries\color{Ink}}
  {\thesubsubsection}{0.6em}{}

\titlespacing*{\section}{0pt}{0.95em}{0.35em}
\titlespacing*{\subsection}{0pt}{0.65em}{0.25em}
\titlespacing*{\subsubsection}{0pt}{0.55em}{0.2em}

% ---------- Table engine ----------
\usepackage{tabularray}
\UseTblrLibrary{booktabs}

% ---------- tabularray: GLOBAL table treatment ----------
\SetTblrInner{
  cells = {fg=black, bg=white, valign=t},
}

% ---------- Header helpers ----------
\newcommand{\Hdr}[1]{\shortstack[c]{\strut #1\strut}}

% ---------- Lines ----------
\newcommand{\TitleLine}{\textcolor{StageBlue}{\rule{0.98\linewidth}{0.95pt}}}
\newcommand{\SubTitleLine}{\textcolor{RuleGray}{\rule{0.98\linewidth}{0.65pt}}}

% ---------- Continued on next page ----------
\newcommand{\ContinuedOnNextPageInline}{%
  \par
  \vfill
  \noindent\textcolor{RuleGray}{\rule{\linewidth}{1pt}}\vspace{-2pt}
  \noindent\hfill\textit{Continued on next page}\par\vspace{-10pt}
  \noindent\textcolor{RuleGray}{\rule{\linewidth}{1pt}}
  \clearpage
}

% ---------- PDF hyperlinks ----------
\usepackage[hidelinks]{hyperref}
\hypersetup{
  colorlinks=false,
  pdfborder={0 0 0},
  linktoc=all,
  pdftitle={\DocTitle},
  pdfauthor={\ProgramName}
}

% =========================================================
\begin{document}

\section*{Stage 2.A — Domain Taxonomy}
\phantomsection
\addcontentsline{toc}{section}{Stage 2.A — Domain Taxonomy}

\subsection*{2.A.1 Taxonomy Principles and Rationale}
\phantomsection
\addcontentsline{toc}{subsection}{2.A.1 Taxonomy Principles and Rationale}

\begin{itemize}
\item Domains exist to lock a controlled measurement vocabulary so evidence remains comparable across 182 weeks and cannot be reinterpreted by convenience, motivation spikes, or documentation drift.
\item Domain tags are governance controls. They are not motivational labels, not “training styles,” and not narrative framing.
\item Domain tags are crosswalk keys. Stage 8 auditing depends on stable domain tagging to reconcile standards, tests, phases, and resources without ambiguity.
\item Domain tags prevent interpretation drift by defining what a measurement construct IS and IS NOT (for example, low-impact aerobic conditioning \textbf{MUST} NOT be treated as running performance or running impact tolerance).
\item Domain tags prevent double-counting by enforcing a single owner construct per metric/test (exactly one Primary Domain Tag).
\item Domain tags protect safety governance by preventing high-consequence constructs (impact running, load carriage, aquatic governance) from being relabeled as “recovery” or “conditioning” to downplay controls.
\item Domain tags protect validity and admissibility by making misclassification visible and correctable; validity tagging (\textbf{VALID} / \textbf{INVALID}) is separate from domain tagging, but both must be consistent for auditability.
\item Domain tags are required fields in metrics (Stage 2.B), protocol cards (Stage 2.C), thresholds (Stage 2.D), and logging conventions (Stage 2.E). Missing or non-canonical tags break crosswalk integrity.
\item Misclassification is an evidence integrity risk. If a metric/test is tagged to the wrong domain, downstream gates, thresholds, and crosswalks become non-auditable and may enable unsafe progression decisions.
\end{itemize}

\subsubsection*{2.A.1.1 Definitions and Terms}
\phantomsection

\begin{itemize}
\item Domain Tag: the canonical domain label assigned to a metric/test to define “what construct this evidence is.” Domain Tags are stable join keys.
\item Domain \textbf{ID}: the stage-scoped identifier for the domain entry in this taxonomy table (DOM-2A-\#\#\#). Domain IDs are stable join keys.
\item Domain Name: plain-language label for the domain construct (informational; does not replace the tag).
\item Metric: a named measurement construct defined in Stage 2.B with an \textbf{ID} (\textbf{MET-2B}-\#\#\#), unit, directionality, and protocol card references.
\item Test: a measurement event executed under a protocol card (Stage 2.C) that produces a recorded result for one or more metrics.
\item Primary Domain Tag: the single construct owner domain tag required for every metric/test. Primary is mandatory.
\item Secondary Domain Tag: an optional modifier domain tag used only to improve crosswalk interpretation (for example, performance constrained by durability flags) without transferring ownership or allowing double credit.
\item Subdomain Directory: internal descriptors that refine interpretation within a domain (e.g., \textbf{RUN-PERF} vs \textbf{RUN-TOL}). Subdomains are not primary tags and must not replace the primary domain tag.
\item Evidence: logged results plus required context fields that permit admissibility determination (\textbf{VALID}/\textbf{INVALID}) under Stage 1.F and Stage 2.E.
\end{itemize}

\newpage

\subsubsection*{2.A.1.2 Non-Negotiables}
\phantomsection

\begin{itemize}
\item Exactly one Primary Domain Tag per metric/test.
\item Primary Domain Tag must be selected from the Canonical Domain List Locked in 2.A.2.
\item Secondary Domain Tag is optional and limited to one (see 2.A.3).
\item Domain tagging does not change validity. Domain tagging classifies the construct; validity tags classify admissibility.
\item Domain tagging must not be altered for optics, convenience, or to bypass governance controls.
\item Domain tags must be applied consistently across:
  \begin{itemize}
  \item metric inventory (2.B),
  \item protocol cards (2.C),
  \item thresholds (2.D),
  \item logging conventions (2.E),
  \item phase references and crosswalks (Stages 3–5 and Stage 8).
  \end{itemize}
\end{itemize}

\clearpage
\begin{landscape}

\subsection*{2.A.2 Canonical Domain List Locked}
\phantomsection
\addcontentsline{toc}{subsection}{2.A.2 Canonical Domain List: Locked}

\begingroup
\scriptsize
\setlength{\tabcolsep}{2pt}
\renewcommand{\arraystretch}{1.12}

\begin{longtblr}{
  width=\linewidth,
  rowhead=1,
  colspec = {
    Q[l,wd=0.65in]
    Q[l,wd=0.65in]
    X[l,1.05]
    X[l,1.35]
    X[l,1.70]
    X[l,1.45]
    Q[l,wd=1.35in]
  },
  hlines,
  vlines,
  cells = {hsep=6pt, vsep=6pt, halign=l, valign=t},
  cell{1-Z}{1-7} = {fg=black},
  row{odd} = {bg=white},
  row{even} = {bg=LightGray},
  row{1} = {bg=StageBlue!15, fg=black, font=\bfseries, halign=l, valign=t},
}
\Hdr{Domain \textbf{ID}} &
\Hdr{Domain Tag} &
\Hdr{Domain Name} &
\Hdr{Definition} &
\Hdr{Includes Examples} &
\Hdr{Excludes Examples} &
\Hdr{Primary Consumers Stages} \\
\textbf{DOM-2A-001} & \textbf{RUN} & Running Performance and Running Tolerance (Impact-Managed) & Measures running-specific performance and running-specific tolerance under impact mechanics. \textbf{RUN} evidence is not interchangeable with low-impact conditioning evidence. & Timed 1.5-mile/2-mile/3-mile/5-mile run trials; run interval benchmarks when defined as \textbf{RUN} tests; run tolerance markers where running mechanics are present and recorded as such. & Bike/rower/elliptical time trials; treadmill walking; generic “cardio minutes” without a running construct; ruck time trials; “felt easy” sessions mislabeled as recovery. & Stage 2.B–2.E; Stage 3; Stages 4–5; Stage 7; Stage 8 \\
\textbf{DOM-2A-002} & \textbf{RUCK} & Load Carriage Performance and Load Carriage Tolerance & Measures load carriage performance under a declared ruck construct, including load verification posture and configuration comparability requirements. & Timed ruck trials (distance/time) with documented load; ruck configuration lock notes used to preserve comparability; ruck tolerance markers tied to load carriage exposure. & Running time trials; unloaded hiking/walking without a declared ruck construct; farmer carries or non-ruck carries (tag \textbf{STR}); swim tests. & Stage 2.B–2.E; Stage 3; Stages 4–5; Stage 6; Stage 8 \\
\textbf{DOM-2A-003} & \textbf{SWIM} & Aquatic Performance (Surface) and Governance-Constrained Underwater Constructs & Measures aquatic performance under unit-locked swim constructs and includes governance-constrained underwater/breath-hold constructs where explicitly permitted by safety prerequisites. & Surface 500-yard/500-meter swim time trials (unit locked); swim technique/efficiency markers tied to swim performance; underwater exposure admissibility flags; underwater 25 m performance only when governance prerequisites are met. & Any unsupervised underwater/breath-hold attempt; dry-land conditioning mislabeled as swim; non-aquatic cardio time. & Stage 2.B–2.E; Stage 3; Stages 4–5; Stage 6; Stage 8 \\
\textbf{DOM-2A-004} & \textbf{CAL} & Calisthenics Performance (Strict Standards) & Measures strict bodyweight performance under standardized rep/hold rules where the construct is calisthenics capacity rather than external-load strength. & Push-ups (timed strict); sit-ups (timed standard); pull-ups (strict dead-hang); plank holds; other strict bodyweight standards explicitly adopted as \textbf{CAL} tests. & External-load strength tests (tag \textbf{STR}); generic “core work” volume; low-impact cardio time; ruck/swim/run tests. & Stage 2.B–2.E; Stage 3; Stages 4–5; Stage 8 \\
\textbf{DOM-2A-005} & \textbf{STR} & Strength and External Load Capacity (Non-Ruck) & Measures strength and external-load capacity constructs that are not ruck time trials, including standardized strength proxies and non-ruck loaded carries. & Strength proxy tests (squat/hinge/press/pull patterns); load × reps strength estimates where defined; loaded carry tests that are not executed as ruck constructs (for example, farmer carry distance/time). & Timed ruck trials (tag \textbf{RUCK}); calisthenics strict rep tests (tag \textbf{CAL}); generic work capacity circuits tagged as \textbf{DURABILITY} in governance posture; aerobic machines used for conditioning performance (tag \textbf{AER}). & Stage 2.B–2.E; Stage 3; Stages 4–5; Stage 6; Stage 8 \\
\textbf{DOM-2A-006} & \textbf{BODYCOMP} & Body Composition and Anthropometrics & Measures bodyweight and body composition trend constructs using method-recorded measurement posture for comparability. & Bodyweight (single measurement and/or weekly average); body fat percent (method/device recorded; trend validity requires method consistency); waist/tape anthropometrics; circumference sites when explicitly adopted. & Performance outcomes (run/ruck/swim/calisthenics/strength); recovery readiness markers (sleep/fatigue); pain flags. & Stage 2.B–2.E; Stages 4–5; Stage 8; Stage 10 \\
\textbf{DOM-2A-007} & \textbf{DURABILITY} & Durability, Structural Integrity Flags, and Disqualifier Signals & Owns non-diagnostic durability signals and constraint flags that control exposure eligibility, test admissibility posture, and conservative gate decisions. \textbf{DURABILITY} is not a “PR domain.” & Pain severity + location; gait-altering pain flags; foot hotspot/blister flags; acute injury instability/swelling flags; movement quality screens recorded as pass/flag under non-diagnostic posture. & Timed performance tests; “recovery day” narratives; medical diagnosis or treatment claims; prescription guidance. & Stage 2.E; Stage 1.F (interface); Stages 4–5; Stage 8 \\
\textbf{DOM-2A-008} & \textbf{RECOVERY} & Recovery and Readiness Markers (Decision Support) & Owns readiness and recovery markers used to protect safety, preserve validity, and constrain gate decisions. \textbf{RECOVERY} does not own performance outputs. & Sleep hours; soreness/fatigue markers; session RPE; resting/trend markers when recorded as non-diagnostic trends; compliance-support markers used to assess whether evidence is admissible and repeatable. & Running/rucking/swimming/calisthenics/strength performance outcomes; relabeling performance constructs as recovery to downplay consequence. & Stage 2.E; Stage 1.F (interface); Stages 4–5; Stage 8 \\
\textbf{DOM-2A-009} & \textbf{AER} & Low-Impact Aerobic Conditioning Performance (Non-Run) & Owns non-impact or low-impact aerobic conditioning performance constructs that are not \textbf{RUN}, \textbf{RUCK}, or \textbf{SWIM}. \textbf{AER} evidence \textbf{MUST} NOT be treated as interchangeable proof of \textbf{RUN} tolerance or \textbf{RUN} performance. & Bike time trials; rower/erg time trials; elliptical time trials; step tests; walk-based aerobic checks where the construct is explicitly non-running conditioning performance. & Timed run trials; run tolerance constructs under impact mechanics; ruck time trials; swim tests; labeling \textbf{AER} as \textbf{RUN} to inflate readiness claims. & Stage 2.B–2.E; Stage 3; Stages 4–5; Stage 8 \\
\end{longtblr}

\endgroup

\textbf{AER} exists to prevent non-run, low-impact conditioning constructs from being mislabeled as \textbf{RUN}, \textbf{RUCK}, or \textbf{SWIM}, which would contaminate comparability and enable unsafe impact escalation.

\end{landscape}
\clearpage

\subsubsection*{2.A.2.1 Domain Tag Formatting and Storage Rules}
\phantomsection

\begin{itemize}
\item Domain Tag storage format: uppercase token exactly matching the Canonical Domain Tag column in 2.A.2.
\item Domain Tag must be stored as a single token (no spaces, no slashes, no synonyms).
\item Domain Tag must be used consistently in:
  \begin{itemize}
  \item \textbf{MET-2B}-\#\#\# records,
  \item \textbf{C1}–\textbf{C18} protocol card headers and validity condition sections,
  \item threshold tables and tier labels,
  \item test logs and evidence packets,
  \item crosswalk joins in Stage 8.
  \end{itemize}
\item Domain \textbf{ID} and Domain Tag must remain 1:1 mapped. Domain \textbf{ID} is the stage-scoped row identifier; Domain Tag is the canonical token used everywhere else.
\item Any appearance of “ROCK” is non-canonical and must be treated as an error. The only canonical tag is \textbf{RUCK}. Any artifact containing “ROCK” is Draft — Non-Deployable until corrected.
\end{itemize}

\subsubsection*{2.A.2.2 Primary Construct Ownership Test}
\phantomsection

A metric or test’s Primary Domain Tag is determined by the construct being measured. Apply the first matching rule below:

\begin{enumerate}
\item If the measured construct is a timed run or running exposure where running mechanics and impact tolerance are present → Primary = \textbf{RUN}.
\item Else if the construct is a timed load carriage event with declared load and ruck configuration posture → Primary = \textbf{RUCK}.
\item Else if the construct is an aquatic performance event (surface swim) or a governance-constrained underwater construct under explicit prerequisites → Primary = \textbf{SWIM}.
\item Else if the construct is strict bodyweight calisthenics under standardized rep/hold rules → Primary = \textbf{CAL}.
\item Else if the construct is external load capacity (strength proxies, non-ruck carries) → Primary = \textbf{STR}.
\item Else if the construct is body composition or anthropometrics → Primary = \textbf{BODYCOMP}.
\item Else if the construct is a non-diagnostic durability flag/disqualifier signal (pain, gait, foot hotspots, instability) → Primary = \textbf{DURABILITY}.
\item Else if the construct is recovery/readiness decision support (sleep, soreness/fatigue, RPE, trend markers) → Primary = \textbf{RECOVERY}.
\item Else if the construct is low-impact aerobic conditioning performance explicitly not running, not ruck, not swim → Primary = \textbf{AER}.
\item If none match, the construct is not defined correctly and must not be instantiated as a metric/test until classified.
\end{enumerate}

\subsubsection*{2.A.2.3 Domain-Specific Boundary Tests}
\phantomsection

\textbf{RUN} boundary tests (impact-managed)

\begin{itemize}
\item If the measurement requires running mechanics under impact → \textbf{RUN}.
\item If the measurement is a walk, incline walk, bike, rower, elliptical, step test, or any non-running construct → not \textbf{RUN}; classify as \textbf{AER} (or another domain as applicable).
\item If “running tolerance” is being claimed without running mechanics exposure → invalid claim; classify as \textbf{AER} and/or \textbf{DURABILITY} flags may apply, but \textbf{RUN} must not be used.
\end{itemize}

\textbf{RUCK} boundary tests (load-carriage construct)

\begin{itemize}
\item If the measurement is timed distance under declared carried load with a ruck configuration posture → \textbf{RUCK}.
\item If load is absent, not declared, or not verified, it is not a ruck construct; it may be \textbf{AER-WALK} or \textbf{DURABILITY} or \textbf{RECOVERY} supportive evidence depending on intent, but must not be tagged \textbf{RUCK}.
\end{itemize}

\textbf{SWIM} boundary tests (aquatic construct)

\begin{itemize}
\item If the measurement is in-water surface swim performance under unit-locked distance → \textbf{SWIM}.
\item If underwater/breath-hold is attempted without governance prerequisites → not a valid \textbf{SWIM} test; it is a governance breach and the evidence is inadmissible regardless of performance.
\end{itemize}

\textbf{CAL} boundary tests (strict bodyweight)

\begin{itemize}
\item If the measurement is strict bodyweight repetition/hold performance under standardized calisthenics rules → \textbf{CAL}.
\item If external load is a defining feature (weighted pull-ups, weighted carries, loaded strength tests) → not \textbf{CAL}; classify \textbf{STR} or other applicable domain.
\end{itemize}

\textbf{STR} boundary tests (external load, non-ruck)

\begin{itemize}
\item If the measurement is external load capacity (strength proxies, non-ruck carries) → \textbf{STR}.
\item If the measurement is executed as a ruck construct (load carriage with pack over distance/time) → not \textbf{STR}; classify \textbf{RUCK}.
\end{itemize}

\textbf{BODYCOMP} boundary tests (composition constructs)

\begin{itemize}
\item If the measurement is bodyweight, body fat percent, waist/tape, anthropometrics → \textbf{BODYCOMP}.
\item If the measurement is performance output (time, reps, distance under performance constructs) → not \textbf{BODYCOMP}.
\end{itemize}

\textbf{DURABILITY} boundary tests (flags/disqualifiers)

\begin{itemize}
\item If the measurement is a non-diagnostic constraint/flag (pain, gait alteration, foot hotspot) used to constrain exposure and admissibility → \textbf{DURABILITY}.
\item \textbf{DURABILITY} must not own timed performance outcomes.
\end{itemize}

\textbf{RECOVERY} boundary tests (readiness markers)

\begin{itemize}
\item If the measurement is readiness/recovery marker (sleep, soreness/fatigue, RPE, trend marker) used for decision support → \textbf{RECOVERY}.
\item \textbf{RECOVERY} must not absorb performance outcomes to downplay consequence.
\end{itemize}

\textbf{AER} boundary tests (low-impact aerobic conditioning performance)

\begin{itemize}
\item If the construct is aerobic conditioning performance and is not running, not rucking, not swimming → \textbf{AER}.
\item \textbf{AER} must not be relabeled as \textbf{RUN}. \textbf{AER} evidence is not interchangeable proof of \textbf{RUN} performance or \textbf{RUN} tolerance.
\end{itemize}

\newpage

\subsection*{Subdomain Directory}
\phantomsection
\addcontentsline{toc}{subsection}{2.A.2.4 Subdomain Directory}

\begin{itemize}
\item \textbf{RUN}

  \begin{itemize}
  \item \textbf{RUN-PERF}: Timed run performance constructs (1.5-mile/2-mile/3-mile/5-mile) under recorded conditions.
  \item \textbf{RUN-TOL}: Running tolerance constructs where impact mechanics are present and durability signals determine whether running exposure is permitted.
  \item \textbf{RUN-INT}: Interval benchmark constructs executed as running (fixed distance/rest) and treated as \textbf{RUN} evidence.
  \end{itemize}

\item \textbf{RUCK}

  \begin{itemize}
  \item \textbf{RUCK-TIME}: Timed ruck constructs (distance/time) under declared load and route verification posture.
  \item \textbf{RUCK-LOAD}: Load verification and configuration lock constructs required to preserve ruck test comparability (supporting admissibility).
  \item \textbf{RUCK-TOL}: Load carriage tolerance constructs (feet/shoulder interface constraints) that govern whether ruck exposure escalation is permitted.
  \end{itemize}

\item \textbf{SWIM}

  \begin{itemize}
  \item \textbf{SWIM-SURF}: Surface swim performance constructs with unit-locked distance and recorded pool verification.
  \item \textbf{SWIM-UW-GOV}: Underwater/breath-hold constructs that exist only under governance prerequisites; evidence is inadmissible without certified supervision.
  \item \textbf{SWIM-CONF}: Water confidence constructs (surface-only) that support safe progression without implying underwater permission.
  \end{itemize}

\item \textbf{CAL}

  \begin{itemize}
  \item \textbf{CAL-PUSH}: Push-dominant calisthenics constructs (push-ups) under strict standards.
  \item \textbf{CAL-PULL}: Pull-dominant calisthenics constructs (pull-ups) under strict standards.
  \item \textbf{CAL-CORE}: Trunk endurance constructs (plank; sit-ups) under strict standards.
  \end{itemize}

\item \textbf{STR}

  \begin{itemize}
  \item \textbf{STR-LOWER}: Lower-body strength proxy constructs (squat pattern) under standardized setup and safe loading posture.
  \item \textbf{STR-HINGE}: Hinge/posterior chain strength proxy constructs (deadlift pattern) under standardized setup and safe loading posture.
  \item \textbf{STR-UPPER}: Upper-body strength proxy constructs (press/pull patterns) under standardized setup and safe loading posture.
  \item \textbf{STR-CARRY}: Non-ruck loaded carry constructs where the measurement construct is strength/load capacity rather than ruck performance.
  \end{itemize}

\item \textbf{BODYCOMP}

  \begin{itemize}
  \item \textbf{BODYCOMP-WT}: Bodyweight constructs (single measurement and/or weekly average) with consistent measurement conditions.
  \item \textbf{BODYCOMP-BF}: Body fat percent constructs where method/device and conditions are recorded every time to preserve trend validity.
  \item \textbf{BODYCOMP-TAPE}: Waist and anthropometric tape constructs with consistent sites and technique posture.
  \end{itemize}

\newpage

\item \textbf{DURABILITY}

  \begin{itemize}
  \item \textbf{DURABILITY-PAIN}: Pain severity/location constructs used as safety and progression constraints (non-diagnostic).
  \item \textbf{DURABILITY-GAIT}: Mechanics and gait-altering flags that function as disqualifiers and regression triggers.
  \item \textbf{DURABILITY-FOOT}: Foot hotspot/blister constructs that constrain impact/load exposure and invalidate “push through” behavior.
  \item \textbf{DURABILITY-MOV}: Movement quality / mobility screens recorded as pass/flag controls (non-diagnostic; supports safe exposure gating).
  \item \textbf{DURABILITY-EVENT}: Acute injury/instability and stop-rule event flags recorded as governance controls (non-diagnostic).
  \end{itemize}

\item \textbf{RECOVERY}

  \begin{itemize}
  \item \textbf{RECOVERY-READ}: Sleep, soreness/fatigue, and session RPE readiness constructs used to constrain testing and interpret validity.
  \item \textbf{RECOVERY-TREND}: Resting marker trend constructs recorded as non-diagnostic trends only (where available and permitted).
  \item \textbf{RECOVERY-COMP}: Compliance-support constructs (session validity tag, Cardio Minimum Standard Met, weekly admissible compliance percent) used to determine gate eligibility.
  \end{itemize}

\item \textbf{AER}

  \begin{itemize}
  \item \textbf{AER-WALK}: Walk-based aerobic checks (including treadmill walking) where the construct is explicitly non-running conditioning performance.
  \item \textbf{AER-BIKE}: Bike-based aerobic performance constructs (time trials, fixed-duration output checks) not treated as \textbf{RUN} evidence.
  \item \textbf{AER-ROW}: Rower/erg aerobic performance constructs not treated as \textbf{RUN} evidence.
  \item \textbf{AER-ELLIP}: Elliptical/cross-trainer aerobic performance constructs not treated as \textbf{RUN} evidence.
  \item \textbf{AER-STEP}: Step-test constructs used as low-impact conditioning checks with recorded setup for comparability.
  \end{itemize}
\end{itemize}

\newpage

\subsection*{2.A.3 Tagging Rules: Primary and Secondary}
\phantomsection
\addcontentsline{toc}{subsection}{2.A.3 Tagging Rules: Primary and Secondary}

Primary tagging rules (mandatory):

\begin{itemize}
\item Every metric and every test \textbf{MUST} have exactly one Primary Domain Tag selected from the canonical domain tags in 2.A.2.
\item The Primary Domain Tag identifies the measurement construct owner. Owner means: the construct is interpreted, thresholded, and crosswalked under that domain.
\item A metric/test MAY carry one Secondary Domain Tag (optional) only when it materially improves crosswalk interpretation (for example, a performance construct that is gated by durability signals).
\item If a Secondary Domain Tag is used, it \textbf{MUST} be treated as a modifier only. It \textbf{MUST} NOT be treated as a second owner, \textbf{MUST} NOT be used to double-count evidence, and \textbf{MUST} NOT be used to “claim credit” in two domains.
\item Metrics and tests \textbf{MUST} NOT be re-tagged to avoid scrutiny, controls, or gate disqualification. Retagging for optics is a governance violation and an evidence integrity breach.
\end{itemize}

Correct tagging examples (minimum set; primary always required; secondary shown only when justified):

\begin{itemize}
\item 1.5-mile run time trial → Primary: \textbf{RUN}; Secondary: \textbf{DURABILITY} (optional when the result is used in a gate where gait flags are disqualifiers).
\item Ruck time trial (measured route, load documented) → Primary: \textbf{RUCK}; Secondary: \textbf{DURABILITY} (optional when foot/hotspot constraints are gate-relevant).
\item Push-ups (2-minute timed strict) → Primary: \textbf{CAL}.
\item Sit-ups (2-minute timed standard) -> Primary: \textbf{CAL}.
\item Pull-ups (strict dead-hang) → Primary: \textbf{CAL}.
\item Plank hold (forearm) → Primary: \textbf{CAL}.
\item Strength proxy (squat/hinge/press/pull pattern) → Primary: \textbf{STR}.
\item Bodyweight (weekly average) / body fat percent (method recorded) / waist tape → Primary: \textbf{BODYCOMP}.
\item Sleep hours / soreness-fatigue / session RPE → Primary: \textbf{RECOVERY}.
\item Gait-altering pain flag / foot hotspot flag → Primary: \textbf{DURABILITY}.
\item Rower or bike time trial / elliptical time trial / step test / walk-based aerobic check (non-run construct) → Primary: \textbf{AER}.
\end{itemize}

Prohibited tagging behaviors (non-negotiable):

\begin{itemize}
\item Labeling a \textbf{RUN} construct as \textbf{RECOVERY} to downplay impact, consequence, or gate disqualification risk.
\item Labeling \textbf{AER} constructs as \textbf{RUN}, including “because it felt hard” or “because it improved aerobic fitness.”
\item Assigning multiple Primary Domain Tags to a single metric/test (double ownership is prohibited).
\item Using a Secondary Domain Tag to obscure ownership, inflate progress, or create double-counted gate credit.
\item Re-tagging after the fact to make a week “look compliant” or to avoid documenting that evidence is not admissible.
\item Using non-canonical domain tags, synonyms, or “custom” labels that fracture crosswalk integrity.
\end{itemize}

\subsubsection*{2.A.3.1 Secondary Domain Tag Allowance}
\phantomsection

Secondary tag is permitted only when it adds crosswalk interpretability without changing ownership.

Permitted patterns (examples):

\begin{itemize}
\item \textbf{RUN} (Primary) + \textbf{DURABILITY} (Secondary) when gait/foot flags materially constrain admissibility or gate interpretation.
\item \textbf{RUCK} (Primary) + \textbf{DURABILITY} (Secondary) when foot hotspots or mechanics-altering pain determine whether ruck evidence can be credited.
\item \textbf{SWIM} (Primary) + \textbf{DURABILITY} (Secondary) only when the metric/test includes explicit durability disqualifiers (e.g., shoulder pain flags) without redefining the primary construct.
\item \textbf{RECOVERY} is generally not a valid Secondary tag for performance tests. Recovery markers influence readiness but do not co-own performance constructs. If needed, reference recovery constraints in validity conditions rather than applying \textbf{RECOVERY} as a secondary performance tag.
\end{itemize}

Forbidden secondary usage:

\begin{itemize}
\item Using \textbf{AER} as Secondary on \textbf{RUN} tests to claim “aerobic transfer.”
\item Using \textbf{RECOVERY} as a secondary tag on \textbf{RUN} or \textbf{RUCK} constructs to hide consequence.
\item Using \textbf{BODYCOMP} as Secondary on performance constructs to imply performance equals composition change.
\end{itemize}

\subsubsection*{2.A.3.2 Composite Evidence and Multi-Metric Test}
\phantomsection

If a single test event yields multiple metrics:

\begin{itemize}
\item Each metric retains its own Primary Domain Tag.
\item The test event itself must have one Primary Domain Tag that matches the primary intent of the event (protocol card determines intent).
\item Supporting fields (load confirmation, footwear record, environment record) do not change domain ownership; they support admissibility.
\end{itemize}

Example handling:

\begin{itemize}
\item A ruck test event producing:
  \begin{itemize}
  \item ruck time metric (Primary \textbf{RUCK}),
  \item load confirmation metric (Primary \textbf{RUCK}; subdomain \textbf{RUCK-LOAD}),
  \item foot hotspot flag (Primary \textbf{DURABILITY}),
  \end{itemize}
  is not “multi-primary.” Each metric is singly owned.
\end{itemize}

\subsubsection*{2.A.3.3 Retagging and Correction Rules}
\phantomsection

\begin{itemize}
\item If a metric/test is found to be misclassified, the correction must be explicit and traceable:
  \begin{itemize}
  \item original domain tag,
  \item corrected domain tag,
  \item rationale referencing taxonomy boundary definitions,
  \item date of correction,
  \item owner and reviewer acknowledgment.
  \end{itemize}
\item Misclassification correction triggers a crosswalk impact review:
  \begin{itemize}
  \item check thresholds, phase references, and any gate interpretations that consumed the mis-tagged evidence,
  \item determine whether any gate decisions were contaminated and require rerun under admissible evidence.
  \end{itemize}
\item Misclassification is not a “cosmetic fix.” Until corrected and impact-checked, affected artifacts are Draft — Non-Deployable.
\end{itemize}

\subsection*{2.A.4 Domain Boundary Notes High Risk Clarifications}
\phantomsection
\addcontentsline{toc}{subsection}{2.A.4 Domain Boundary Notes High Risk Clarifications}

\begin{itemize}
\item \textbf{RUN} boundary (impact governance)

  \begin{itemize}
  \item \textbf{RUN} owns running-specific constructs under impact mechanics.
  \item Low-impact conditioning improvements (\textbf{AER}) \textbf{MUST} NOT be interpreted as \textbf{RUN} tolerance proof or \textbf{RUN} performance proof.
  \item If running mechanics are absent (no running construct), the domain \textbf{MUST} NOT be \textbf{RUN}.
  \item If a test is “treadmill walking,” it is \textbf{AER-WALK}, not \textbf{RUN}, even if performed at high effort.
  \end{itemize}

\item \textbf{SWIM} boundary (surface-only unless governed)

  \begin{itemize}
  \item \textbf{SWIM} includes surface swim performance and governance-constrained underwater constructs.
  \item Underwater or breath-hold exposure is governed and \textbf{MUST} NOT occur unsupervised under any circumstances.
  \item Any underwater/breath-hold attempt without certified supervision is a safety violation and the resulting evidence is not admissible for readiness claims regardless of performance.
  \item Surface swim performance does not imply underwater permission or underwater readiness.
  \end{itemize}

\item \textbf{AER} boundary (non-impact conditioning performance)

  \begin{itemize}
  \item \textbf{AER} owns low-impact aerobic conditioning performance constructs (bike/row/elliptical/step/walk-based checks) that are not \textbf{RUN}, \textbf{RUCK}, or \textbf{SWIM}.
  \item \textbf{AER} evidence is not a substitute label for \textbf{RUN} readiness and \textbf{MUST} NOT be used to authorize impact escalation.
  \item \textbf{AER} is permitted specifically to prevent mislabeling and evidence contamination.
  \item \textbf{AER} may support general conditioning readiness posture but cannot be used to claim \textbf{RUN} test thresholds were met.
  \end{itemize}

\item \textbf{DURABILITY} boundary (flags and constraints; not a PR domain)

  \begin{itemize}
  \item \textbf{DURABILITY} owns non-diagnostic flags, constraints, and disqualifier signals (pain, gait change, hotspots, instability, movement screens).
  \item \textbf{DURABILITY} is not a performance PR domain and \textbf{MUST} NOT be used to “win” numbers; it exists to constrain risk and protect admissibility and safe progression.
  \item \textbf{DURABILITY} constructs \textbf{MUST} be treated as decision constraints that can force \textbf{HOLD}, \textbf{REGRESS}, or \textbf{MEDICAL} \textbf{HOLD} posture when triggered (per upstream governance).
  \item \textbf{DURABILITY} signals are not diagnoses and must not be described as diagnoses.
  \end{itemize}

\item \textbf{RECOVERY} boundary (decision support for validity and risk; not performance)

  \begin{itemize}
  \item \textbf{RECOVERY} owns readiness and recovery markers (sleep, fatigue, soreness, RPE, trend markers where applicable) and compliance-support markers used for gate eligibility.
  \item \textbf{RECOVERY} \textbf{MUST} NOT absorb performance outcomes "because it was easy."
  \item \textbf{RECOVERY} constructs inform whether performance evidence can be safely expressed and whether testing is admissible; they do not replace performance domain evidence.
  \item \textbf{RECOVERY} does not provide gate credit; it constrains whether gate credit can be safely pursued.
  \end{itemize}
\end{itemize}

\newpage

\subsection*{2.A.5 Taxonomy Change Control Requirements}
\phantomsection
\addcontentsline{toc}{subsection}{2.A.5 Taxonomy Change Control Requirements}

\begin{itemize}
\item Domain IDs and Domain Tags are frozen once published.
\item Domain IDs are stable crosswalk keys and \textbf{MUST} NOT be reused, repurposed, or reassigned.
\item Any change to domains, tags, or IDs \textbf{MUST} use \textbf{PATCH}-2A-\#\#\#.
\item Any taxonomy change triggers downstream revalidation at minimum across Stage 2.B through Stage 2.E and requires a crosswalk impact check (Stage 8 impact posture).
\item No silent edits. If taxonomy meaning changes, the change \textbf{MUST} be explicit, logged, and traceable.
\end{itemize}

\subsubsection*{2.A.5.1 Patch \textbf{ID} Convention}
\phantomsection

\begin{itemize}
\item Taxonomy patches for this deliverable use: \textbf{PATCH-2A-001}, \textbf{PATCH-2A-002}, … sequential with no gaps unless retired.
\item Patch record must include:
  \begin{itemize}
  \item Patch \textbf{ID},
  \item date,
  \item change description,
  \item rationale,
  \item affected domains/tags,
  \item downstream artifacts impacted,
  \item revalidation scope,
  \item owner.
  \end{itemize}
\end{itemize}

\subsubsection*{2.A.5.2 Change-Type Table}
\phantomsection

\begingroup
\scriptsize
\setlength{\tabcolsep}{2pt}
\renewcommand{\arraystretch}{1.12}

\begin{longtblr}{
  width=\linewidth,
  rowhead=1,
  colspec = {
    X[l,1.15]
    X[l,1.55]
    Q[l,wd=1.10in]
    X[l,2.35]
  },
  hlines,
  vlines,
  cells = {hsep=6pt, vsep=6pt, halign=l, valign=t},
  cell{1-Z}{1-4} = {fg=black},
  row{odd} = {bg=white},
  row{even} = {bg=LightGray},
  row{1} = {bg=StageBlue!15, fg=black, font=\bfseries, halign=l, valign=t},
}
\Hdr{Change Type} & \Hdr{Examples} & \Hdr{Requires Patch} & \Hdr{Revalidation Scope} \\
Clarification (no meaning change) & Add examples; tighten wording without changing boundaries; improve definition clarity & Yes (\textbf{PATCH}-2A-\#\#\#) & Stage 2.B–2.E compatibility check + crosswalk impact check to confirm no downstream re-tagging is required \\
Boundary change (meaning change) & Move a construct between domains; redefine \textbf{RUN} vs \textbf{AER} boundary; redefine what \textbf{DURABILITY} owns & Yes (\textbf{PATCH}-2A-\#\#\#) & Stage 2.B–2.E revalidation required + crosswalk impact check; affected downstream artifacts default Draft — Non-Deployable until revalidated \\
Domain addition & Introduce a new domain tag because an evidence construct cannot be classified without ambiguity & Yes (\textbf{PATCH}-2A-\#\#\#) & Full Stage 2.B–2.E revalidation + crosswalk impact check; downstream artifacts remain Draft — Non-Deployable until updated to use the new domain \\
Domain retirement & Retire a domain \textbf{ID}/tag due to consolidation; preserve non-reuse rule & Yes (\textbf{PATCH}-2A-\#\#\#) & Stage 2.B–2.E revalidation + crosswalk impact check; retired \textbf{ID} remains reserved and \textbf{MUST} NOT be reused \\
\textbf{ID} correction (administrative but crosswalk-critical) & Fix an \textbf{ID} formatting defect; correct sequencing record while preserving non-reuse & Yes (\textbf{PATCH}-2A-\#\#\#) & Full crosswalk impact check required; Stage 2.B–2.E reviewed for join-key integrity; any affected artifacts default Draft — Non-Deployable until corrected \\
\end{longtblr}

\endgroup

\subsection*{2.A.6 Conflict-Resolution Rule}
\phantomsection
\addcontentsline{toc}{subsection}{2.A.6 Conflict-Resolution Rule}

If any downstream artifact uses non-canonical domain tags, violates the one-primary rule, violates the one-secondary limit, or ignores boundary definitions, that artifact is Draft — Non-Deployable until corrected and re-versioned.

\end{document}